\documentclass{scrartcl}
\usepackage[utf8]{inputenc}
\usepackage{graphicx}
\usepackage{amsmath}
\usepackage{float}
\usepackage{siunitx}
\sisetup{
  per-mode=symbol,
  separate-uncertainty
}
\usepackage{hyperref}
\hypersetup{
    colorlinks=true,
    linkcolor=blue,
    filecolor=magenta,      
    urlcolor=cyan,
    pdfpagemode=FullScreen,
    }
\usepackage{authblk}

\title{Title}

\author[1]{Robin Kühnberger}
\author[2]{Felix Naraschewski}

\affil[1]{University Stuttgart}
\affil[2]{University Stuttgart}

\date{\today}

\begin{document}
\maketitle
\tableofcontents

\section{section0}

\subsection*{Floats}
Floats wie figure und table werden von \LaTeX{} automatisch so platziert, dass ein guter Seitenumbruch und ein sauberer Textfluss entstehen. Sie stehen daher nicht zwingend an der Stelle im Code.
\subsubsection*{Positionsparameter}
\begin{itemize}
    \item \textbf{[h]} – here: möglichst an der aktuellen Stelle
    \item \textbf{[t]} – top: am Seitenanfang
    \item \textbf{[b]} – bottom: am Seitenende
    \item \textbf{[p]} – page: auf einer eigenen Float-Seite
    \item \textbf{[H]} – genau hier (erfordert \texttt{\textbackslash usepackage\{float\}})
\end{itemize}


\subsection{section1}
\subsubsection{section2}

\begin{figure}[H]
	\centering
	\includegraphics[width=0.6\textwidth]{starstek.png}
	\caption{United federation of planets - Starfleet Command Logo}
	\label{img:starfleet_command}
\end{figure}

\begin{equation}
\begin{aligned}
    R_{\mathrm{\mu \nu}} - \frac{1}{2}g_{\mathrm{\mu \nu}}=\frac{8 \pi G}{c^{4}}T_{\mathrm{\mu \nu}} \\
    \sum_{i=1}^{n} i = \frac{n(n+1)}{2}
\label{eq:gleichung0}
\end{aligned}
\end{equation}

\begin{table}[!ht]
    \caption{Bekannte Raumschiffe aus Star Trek.}
    \centering
    \begin{tabular}{|l||l|l|}
        \hline
        \textbf{Raumschiff}         & \textbf{Maximaler-Warp-Faktor} & \textbf{Kapitän}                 \\
        \hline
        USS Enterprise (NCC-1701)   & Warp 9.6                       & James T. Kirk                    \\
        \hline
        USS Enterprise (NCC-1701-D) & Warp 9.8                       & Jean-Luc Picard                  \\
        \hline
        USS Voyager (NCC-74656)     & Warp 9.975                     & Kathryn Janeway                  \\
        \hline
        USS Defiant (NX-74205)      & Warp 9.5                       & Benjamin Sisko                   \\
        \hline
        USS Discovery (NCC-1031)    & Warp 9.9                       & Gabriel Lorca / Christopher Pike \\
        \hline
    \end{tabular}
    \label{tab:starships}
\end{table}
Die Abbildung \ref{img:starfleet_command} zeigt das Logo von Starfleet Command. (They really exist!)\\
Raumschiffe sind toll \ref{tab:starships}\\
Die Gleichungen \eqref{eq:gleichung0} sind zufällig ausgewählt.\\

Die Lichtgeschwindigkeit \cite{einstein1966einstein} beträgt \(\num{2.99792458e8}\,\si{\meter\per\second}\),
die Elementarladung \cite{schopper1935neubestimmung} hat den Wert \(\num{1.602176634e-19}\,\si{\coulomb}\)
und die Planck-Konstante \cite{steiner2012history} ist \(\num{6.62607015e-34}\,\si{\joule\second}\).

\section*{Quellen}
\bibliographystyle{plain}
\bibliography{bibliography}

\listoffigures
\listoftables

\end{document}